\chapter{Global cylinder control: arbitrary words, whole crossing tails,
and retained singleton defects}
{\small\sffamily Source: \nolinkurl{repository/collatz_reconstruction/research_program/global_cylinder_control_20260916/note.md}.\par}


16 September 2026. Additive Collatz workbench continuation, read at
\nolinkurl{9afcbe2dc9bace77c1ae705b1cd128da22fe8f8b}, PR \#2. The
construction uses the original +1 map and the published integral
first-jet receiver. It does not import an unproved estimate from the
Zeta programme. The preceding 33-file local arithmetic/shadow package is
not an input. This note makes no priority claim and no assertion of
global Collatz convergence.

The main arithmetic quantifiers run over \textbf{all finite exponent
words}, not only the earlier P/Q languages, repeated cores or a bounded
switch count. The price is explicit: one original source can remain at
each contracting cylinder, and at most one original source remains
across the entire unbounded next-exponent crossing tail of a
pre-crossing prefix. Those sources are retained rather than declared
zero. A separate finite proof certificate discharges the first-crossing
exception through 18 odd returns.

\section{1. Original objects and source maps}

Let X be the positive odd integers and

\SourceMath{T(n)=(3n+1)/2^{a(n)},\qquad a(n)=\nu_2(3n+1).}{}

For an ordered nonempty word p=(a\_1,\ldots,a\_m), a\_i\textgreater=1,
retain

\SourceMath{A_j=\sum_{i=1}^ja_i,\quad A_0=0,\quad
 U=2^{A_m},\quad L=3^m,\quad D=U-L,}{}
\SourceMath{C=C(p)=\sum_{j=0}^{m-1}3^{m-1-j}2^{A_j},\qquad
 f_p(x)=(Lx+C)/U.}{1}

The branch order is f\_p=g\_(a\_m)\ldots g\_(a\_1),
g\_a(x)=(3x+1)/2\^{}a. The complete original positive source is

\SourceMath{\mathcal C_p=\{\rho_p+2Ut:t\in\mathbb Z_{\ge0}\},\quad
 \rho_p\equiv(U-C)L^{-1}\pmod{2U},\quad0<\rho_p<2U.}{2}

Its original image is y\_p+2Lt, y\_p=(L rho\_p+C)/U, and both t inverses
are given by the displayed source or target difference divided by its
original step. In particular the target step is not changed to the
source step. The positive representative is odd: U is even and C is odd.

For completeness, the terminal-odd congruence (2) gives all intermediate
oddness conditions. Backwards from an odd last value, the recurrence
\nolinkurl{3x_j+1=2^(a_(j+1))x_(j+1),} together with the prefix formula,
shows each previous numerator divisible by its power of two and every
previous x\_j odd. Equivalently, modulo 2\^{}(A\_j+1), the composite
numerator has its first j terms plus a last odd contribution 2\^{}A\_j.
Since every a\_i\textgreater=1, terms past that contribution vanish.
Dividing the resulting congruence by 2\^{}A\_j gives odd x\_j. More
explicitly, for j\textless m the complete numerator reduces to
3\^{}(m-j)*(3\^{}j n+C\_j)+3\^{}(m-j-1)*2\^{}A\_j modulo 2\^{}(A\_j+1).
The terminal congruence makes this zero; division by the odd coefficient
and inversion of 3 give 3\^{}j n+C\_j congruent to 2\^{}A\_j modulo
2\^{}(A\_j+1). The exact recurrence then forces the original valuation
a\_(j+1). Positive forward values follow directly from g\_a.

An explicit append map, also used by the checker, starts from n=rho+2Ut
and x=y+2Lt. To append exponent a, put

\SourceMath{k\equiv\left(2^{a-1}-(3y+1)/2\right)(3L)^{-1}\pmod{2^a},
 \quad0\le k<2^a.}{3}

Then the new anchor is rho+2Uk, the new target is (3y+1+6Lk)/2\^{}a, and
the new coefficients are L'=3L, U'=2\^{}a U, C'=3C+U. Substituting
t=k+2\^{}a v proves the whole new progression and its inverse. Empty
prefix data are \nolinkurl{(rho,y,L,U,C)=(1,1,1,1,0);} its source is all
positive odds.

\section{2. A proved global separation of powers}

\subsection{Theorem 1. All-length separation}

For every integer m\textgreater=1 and every integer A with
2\^{}A\textgreater3\^{}m,

\SourceMath{\boxed{\quad 2^A-3^m>(3/2)^m-1.\quad}}{4}

This is a computer-assisted theorem with one stated external input,
Matveev's established rational-number logarithm bound {[}MAT{]}. The
complete finite part is generated and checked in
\nolinkurl{power_gap_certificate.json}. It consists of rational interval
inequalities and 127 integer rows, not a large Collatz-start
computation. The general argument follows.

Take A\_0=floor(m log(3)/log(2))+1; then A\_0\textless=2m. Increasing A
increases the left side of (4), so it suffices to use A\_0. A failure
gives

\SourceMath{0<\Lambda=2^{A_0}3^{-m}-1<2^{-m}.}{5}

Apply {[}MAT, Theorem 2.1{]} with its two positive rational numbers 2,3,
integer exponents A\_0,-m, and B=2m. Their heights are log 2, log 3. Its
coefficient is dominated by

\SourceMath{(7/5)30^5\,23\,(7/10)(11/10)<10^9,}{6}

because 2\^{}(9/2)\textless23, log 2\textless7/10 and log
3\textless11/10. Thus

\SourceMath{-10^9(1+\log(2m))<\log\Lambda<-m\log2<-(2/3)m.}{7}

At m=10\^{}12, use 2*10\^{}12\textless2\^{}41 and log2\textless7/10 to
get

\SourceMath{(2/3)10^{12}>10^9(1+41\cdot7/10).}{}

The difference (2/3)m-10\^{}9(1+log(2m)) has derivative
2/3-10\^{}9/m\textgreater0 for m\textgreater=10\^{}12. Hence a failure
has m\textless10\^{}12.

All logarithm bounds are certified by rational series. For x=2 or 3,
z=(x-1)/(x+1), use

\SourceMath{\log x=2\sum_{j=0}^{N-1}\frac{z^{2j+1}}{2j+1}+R_N,
 \quad0<R_N<\frac{2z^{2N+1}}{(2N+1)(1-z^2)}.}{8}

The remainder bound follows by replacing all later positive denominators
by 2N+1 and summing a geometric series. With N=160, exact rational
arithmetic gives, for alpha=log3/log2,

\SourceMath{\frac{83130157078217}{52449289519716}
 <\alpha<\alpha+2^{-128}
 <\frac{350861368503572}{221368876767703}.}{9}

The outer determinant is one. A rational x/y strictly between p/q and
r/s with rq-ps=1 has y\textgreater=q+s: the positive integers k=xq-py
and j=ry-sx satisfy y=sk+qj. The identity does not require reducing x/y.
Here \nolinkurl{q+s=273818166287419>10^12.}

For m\textgreater=128, (5) gives

\SourceMath{0<A_0/m-\alpha
 =\log(1+\Lambda)/(m\log2)
 <\frac{3}{2m}2^{-m}<2^{-128}.}{10}

This contradicts (9), since its original denominator m is below
10\^{}12. It remains to check 1\textless=m\textless=127. For each m the
checker takes A\_0=bit\_length(3\^{}m) and verifies the strict integer
margin

\SourceMath{((2^{A_0}-3^m)+1)2^m-3^m>0.}{11}

All 127 rows pass and are retained in full. These cases exhaust (4). No
floating-point evaluation or guessed decimal bound occurs. QED.

\section{3. Every contracting word has at most one non-descending
source}

The coefficient-only comparison below concerns the formal affine maps at
x=0. It does not treat zero as an actual positive odd Collatz source.
For x\textgreater=0,

\SourceMath{g_a(x)=2^{1-a}g_1(x)\le g_1(x),\qquad g_1(x)=(3x+1)/2.}{}

Monotonicity, applied at the same input at each step, gives

\SourceMath{0<C/U=f_p(0)\le2^{1-a_m}\big((3/2)^m-1\big).}{12}

The last factor 2\^{}(1-a\_m) is kept separately: bound the preceding
m-1 maps by g\_1\^{}(m-1), then use g\_(a\_m)=2\^{}(1-a\_m)g\_1. This
proves (12) without identifying different ordered affine maps or
dropping C.

\subsection{Theorem 2. All-word canonical-source isolation}

For every original word p with D\textgreater0,

\SourceMath{\boxed{\quad C/D<2^{A_m-a_m+1}=2^{A_{m-1}+1}.\quad}}{13}

The set of its positive sources with f\_p(n)\textgreater=n is exactly

\SourceMath{\boxed{
 \{\rho_p\}\quad\text{when }C\ge D\rho_p,
 \qquad\varnothing\quad\text{when }C<D\rho_p.
 }}{14}

All other sources in the infinite original cylinder strictly descend.
The candidate in (14) is retained even when its resulting set is empty.

\textbf{Proof.} Combine (4) and (12) to get (13). The exact displacement
f\_p(n)-n=(C-Dn)/U makes every non-descending positive n at most C/D.
Every legal source except rho\_p is at least rho\_p+2U\textgreater2U,
whereas the right side of (13) is at most 2U. Thus only rho\_p is
possible. Testing C-D rho\_p gives exactly (14). On the whole source the
unchanged identity is

\SourceMath{n(t)-f_p(n(t))=\rho_p-y_p+2Dt.}{15}

Its positive slope proves the rest after its first accepted parameter.
For D\textless0, the same formula in n gives f\_p(n)\textgreater n on
every positive source, since C\textgreater0. D=0 is impossible for a
nonempty word. QED.

This theorem quantifies over every finite exponent word, with arbitrary
switches, lengths, exponents and orderings. It does not say the
exceptional canonical source always descends. For example the original
word

\begin{Verbatim}[breaklines=true,breakanywhere=true,fontsize=\scriptsize]
(4,1,1,1,1,2,2,1,2,1,1,2,1,1,1,2,3)
\end{Verbatim}

has m=17,A=27, canonical source 165 and endpoint 167, with
D\textgreater0. It already visited 31 on its first return, so this is
not a first-crossing exception. The source and image of its whole
cylinder are

\SourceMath{165+268435456t\longmapsto167+258280326t.}{16}

Every t\textgreater=1 strictly descends; t=0 is explicitly retained and
has its own verified finite arrival witness. This is an example of the
coefficient- contracting yet non-descending behavior investigated in
{[}RT{]}, not a new claim that such behavior is absent.

\section{4. The entire unbounded next-exponent crossing family}

A pre-crossing prefix q of length l has U\_j\textless3\^{}j at every
nonempty prefix. Keep its data (rho,y,U,L,C) as above. Define

\SourceMath{h=\lfloor\log_2(3L)\rfloor-A_l+1\ge2.}{17}

Thus 2\^{}(h-1)U\textless3L\textless2\^{}hU. The next original exponent
has not crossed exactly for a=1,\ldots,h-1. Every a\textgreater=h is a
first coefficient crossing. Neither a frequency model nor independence
of exponents is used.

Let

\SourceMath{M=2^{h-1},\quad t_0\equiv-((3y+1)/2)(3L)^{-1}\pmod M,
 \quad0\le t_0<M,}{} \SourceMath{n_0=\rho+2Ut_0,\quad
 b_0=(3y+1+6Lt_0)/2^h.}{18}

The full union of all crossing children is the one positive progression

\SourceMath{\boxed{
 n=n_0+U2^h v,\quad
 T^{l+1}(n)=\frac{b_0+3Lv}{2^{\nu_2(b_0+3Lv)}},\quad
 a_{l+1}=h+\nu_2(b_0+3Lv),\quad v\ge0.
 }}{19}

The inverse is v=(n-n\_0)/(U2\^{}h) on the displayed source.
Substitution into 3T\^{}l(n)+1 proves (19) and the full valuation,
including arbitrarily large extra divisions. Congruence (18) is
equivalent to its divisibility by 2\^{}h, so the source is complete. The
low children retain (3) and their own labels. Together, they exhaust
every next exponent and every positive source of q.

\subsection{Theorem 3. One original test controls the whole crossing
alphabet}

The only possible non-descending source in (19) is the OLD canonical
prefix source rho. It belongs to (19) precisely when t\_0=0, and its
actual next exponent is uniquely a*=nu\_2(3y+1). It is an exception
precisely when

\SourceMath{a^*\ge h\quad\text{and}\quad (3y+1)/2^{a^*}\ge\rho.}{20}

Every other source in the full unbounded crossing family has its first
strict descent below its start at return l+1.

\textbf{Proof.} For any admitted child p=qa, (13) says a non-descending
source n is less than 2\^{}(A\_l+1)=2U. Its q-cylinder congruence then
forces n=rho. That source has just one actual next exponent, proving
(20). Earlier prefixes have multiplier greater than or equal to one and
positive forcing, so every proper nonempty endpoint exceeds n.~Thus
final descent is its first strict descent. The empty prefix case retains
n=1 as the actual fixed-loop exception and all other crossing sources
descend. QED.

Each individual family also has an elementary checkable proof once its
original coefficients have been calculated. Its undivided target line is
b\_0+3Lv, its source line is n\_0+U2\^{}h v, and

\SourceMath{U2^h>3L,\qquad n_0+U2^h>b_0+3L.}{21}

The second inequality follows from (13) applied to the formal affine
word qh and the positive source n\_0+U2\^{}h\textgreater2U; the
inequality (13) itself is about all real affine inputs and does not need
terminal oddness. The implementation checks both integer inequalities
directly. They certify every v\textgreater=1 by their positive slope;
v=0 is the single exact valuation check. Thus a delivered individual
family can be verified without executing a transcendence theorem or
enumerating its sources.

\subsection{An original nonrepeating example}

For the single mixed prefix

\begin{Verbatim}[breaklines=true,breakanywhere=true,fontsize=\scriptsize]
q=(1,2,1,1,1,2,1,1,4)=PQ,
\end{Verbatim}

retain \nolinkurl{L=19683,U=16384,C=29839,rho=10907,y=13105,h=2.}
Formula (19) is

\SourceMath{\boxed{
 10907+65536v\longmapsto
 \frac{9829+59049v}{2^{\nu_2(9829+59049v)}},\quad v\ge0.
 }}{22}

This is first strict descent after ten odd returns. The actual last
exponent is 2+nu\_2(9829+59049v). Indeed the endpoint is at most
\nolinkurl{9829+59049v<10907+65536v.} Arbitrarily large last exponents
have positive source representatives because 59049 is odd. No repetition
of PQ and no unspecified weight for v occurs.

\section{5. Whole-system finite audit with infinite source parameters}

The exact symbolic audit starts at the empty prefix. At a prefix of
length l it performs the single test (20), checks (21), and creates ALL
low children a=1,\ldots,h-1. The full high tail (19) is retained as a
complete certificate and is not truncated by exponent size. The exact
number of nodes at each depth is independently checked by the recurrence

\SourceMath{N_0(0)=1,\quad
 N_j(A)=\sum_{B<A}N_{j-1}(B),\quad
 j\le A\le\lfloor\log_2(3^j)\rfloor.}{23}

Every original pre-crossing word satisfies precisely these prefix
bounds. Recurrence (23) and the explicit child ranges prove completeness
of the finite tree. At EVERY node the generated rho is checked against
(2) using an independent modular inverse, the original affine equation
U*y=L*rho+C is checked, and the next valuation is recomputed by repeated
division. Both whole-tail margins (21) and the actual endpoint of the
whole crossing family's first source n\_0 are also checked. Thus this
finite audit supplies its own elementary affine-tail proofs and does not
need the external logarithm theorem as a premise: each v=0 is checked
directly, and the positive affine margin at v=1 proves all
v\textgreater=1.

The complete run for first crossing at return at most 18 visits
1,166,058 prefixes. It retains exactly one exceptional crossing: source
1, word (2), endpoint 1. Therefore every positive source n\textgreater=3
whose first odd-return coefficient crossing is at most 18 has first
actual strict descent at that return. This is unbounded in the starting
integer and the last exponent. Its source-time scope is not a bound on
every orbit's crossing return, and it is not presented as a published
stopping-time record.

The ordered stream can be materialized with \nolinkurl{--ledger}. The
independent \nolinkurl{audit_ledger.py} imports neither the main module
nor the predecessor; it reconstructs the entire depth-first domain and
each original numerator, residue, endpoint, unbounded-tail margin and
canonical-point test. Its successful full replay is included in the
delivery evidence. The uncompressed SHA-256 and per-depth counts are in
verification.json. The executable proof certificate regenerates the
ENTIRE finite domain, not a sample. The independent count recurrence
does not replace the original affine and valuation checks. The
infinite-depth frontier remains unproved.

For example n=27 first crosses at return 37, reaches 23 there, and has
an independently checked 41-return witness to 1. It is outside the
18-return finite theorem but its whole crossing cylinder is still
covered by the all-length reduction and a single original root test.

\section{6. Exact transfer into the retained Split-Zero comparison}

The original absolute first-jet complex has original edge and vertex
modules C\^{}0,C\^{}1, source and target maps J,P, d=J-P, and

\SourceMath{d_\epsilon=d-\epsilon P,\quad \epsilon^2=0.}{}

The receiver is the already published kernel

\SourceMath{\mathfrak B_X=
 \ker\big(H^1(K_\epsilon)\to H^1(K)\big).}{}

Its integral presentation is C\^{}1/(dC\^{}0+P ker d), with the inverse
map to first-order classes {[}z{]} -\textgreater{} {[}epsilon z{]}. The
predecessor proves that a finite integral pair db=0, dc-Pb=V\_n yields
an actual finite path to 1, and that the full obstruction has cyclic Z/m
summands and nonperiodic Z summands. Those definitions and source maps
remain unchanged {[}WB1{]}.

Let B\_p(n) be the sum of the actual original edges in an accepted word
from n to y. Its boundary and its new first-order use are

\SourceMath{dB_p(n)=V_n-V_y,\qquad
 d_\epsilon(\epsilon B_p(n))=\epsilon V_n-\epsilon V_y.}{24}

Thus the all-word descent gives {[}epsilon V\_n{]}={[}epsilon V\_y{]}
with y\textless n. It does NOT by itself assert either class is zero. An
original witness at y is transported by the explicit addition

\SourceMath{(b_y,c_y)\longmapsto(b_y,c_y+B_p(n)).}{25}

Substitution in db and dc-Pb proves the exact boundary at n.~No original
edge or cycle is divided by a period and no integral denominator is
removed. The checker imports the unchanged, Git-blob-verified
firstjet.py and verifies (24), including all intermediates and terminal
vertices. For the examples with complete arrival, its unchanged witness
checker extracts the actual path.

At an equation support label E, (24) maps a declared source module into
the original equations in that support, or enlarges E by exactly the
finitely many edges of the displayed path. Unions of such supports give
the original inclusion diagram, and (24)-(25) commute with its
inclusions. Their supported zero is at that retained E. An empty
exceptional set in (14) or (20) is stored alongside its source label and
integer equation, not substituted for the external absent scalar. The
word index itself is not a definition of intrinsic Split-Zero support.

The family maps are finite-column maps on free modules even over the
infinite progression: each original basis vector has its own finite
path. There is no infinite sum in (24). Both the true singleton
exception at 1 and the non-descending canonical point 165 remain
available to the original operator and subsequent certificates.

\section{7. Relation to the elementary maps, with the different source
moduli}

Let F be the shortened elementary map: F(x)=(3x+1)/2 at odd x and
F(x)=x/2 at even x. The word p expands to the parity word

\SourceMath{w(p)=1\,0^{a_1-1}\,1\,0^{a_2-1}\cdots1\,0^{a_m-1},
 \qquad |w|=A_m.}{26}

The composition has the exact same L,U,C as (1). Grouping at the
original odd vertices is its inverse for paths with odd terminal value.
The parity cylinder alone has modulus 2\^{}A; additionally requiring
terminal oddness selects the original odd-return cylinder modulo
2\^{}(A+1). These moduli are not identified. On their common source,
T\^{}m(n)=F\^{}A(n).

For ANY shortened parity word with k letters and m odd letters, its
original affine constant E satisfies 0\textless=E\textless=
(3/2)\^{}m-1: omitting a division only increases the nonnegative formal
constant. If 2\^{}k\textgreater3\^{}m, (4) gives
E\textless2\^{}k-3\^{}m, so its non-descent threshold is strictly below
2\^{}k. Thus every such parity cylinder also has at most its single
canonical positive representative as an exception. The all-even word has
zero constant and strictly decreases every positive source, including
its least positive source 2\^{}k; the canonical residue zero is not a
positive starting point.

The first shortened coefficient crossing can occur inside the final even
divisions of an odd-return block. Consequently the finite theorem of
Section 5 is stated with odd-return counts and is not silently asserted
as the identical count in Terras' shortened coefficient-stopping
problem. The exact expansion (26), its terminal condition and its clocks
are the comparison used here. The standard elementary map without
immediate odd-step division further expands each odd shortened step into
its original odd step plus one division, with inverse grouping.

\section{8. Global restrictions and the remaining source problem}

This is an all-word source-isolation theorem, not a proof that its
canonical candidates are always harmless. For any nonbase component, its
least original positive integer s has every iterate at least s. A
coefficient-crossing prefix at s is therefore a non-descending case of
(14). More precisely, at its first crossing it must pass test (20) as
the canonical source of the proper prefix. The all-word theorem
eliminates every other point of that entire infinite crossing family. If
there is no coefficient crossing, all original prefix multipliers remain
at least one and the associated itinerary stays in the infinite
pre-crossing tree. Both alternatives remain represented.

The same exact inequalities are useful for finite returns. Every actual
positive cycle has 2\^{}A\textgreater3\^{}m, by multiplying the
unchanged cycle equations. At every rotation its starting vertex must be
the canonical representative in (2) and satisfy (13). Its integral
forcing condition is still D\textbar C; no canonical-root test replaces
the actual cyclic solution and valuations.

At the first coefficient crossing m, all proper-prefix
2\^{}A\_j\textless=3\^{}j give C\textless=m*3\^{}(m-1). Hence a
non-descending original endpoint n+2k satisfies

\SourceMath{0\le k<m/6.}{27}

Indeed (f\_p(n)-n)\textless C/U\textless m/3 because U\textgreater3\^{}m
and Dn\textgreater0. For m=1 the same strict upper bound applies to its
equality at 1. This is an additional all-length exact small-overshoot
restriction, not a claim that every allowed k occurs.

There is no fixed finite bound on the first crossing for all sources:
n\_N=2\^{}(N+1)-1 has its first N exponents equal to one, with
T\^{}j(n\_N)=3\^{}j2\^{}(N+1-j)-1 for 0\textless=j\textless=N. All those
prefix coefficients exceed one. The sources vary with N. This explicitly
prevents promotion of the finite tree audit to one uniform clock bound.

The next arithmetic calculation is therefore sharply specified: extend
source proofs on the actual canonical crossing points and control the
ordinary-positive-integer support of the no-crossing infinite branches.
The canonical residue towers have their original increasing
representatives; boundedness gives eventual constancy, not an automatic
contradiction. The current argument does not show their positive
intersections empty. A finite failure to cross is not an infinite-orbit
certificate. The workbench still retains both periodic and nonperiodic
comparison-kernel summands.

\section{9. Verification and source discipline}

\nolinkurl{power_gap_certificate.json} contains the exact finite
arithmetic component of Theorem 1. \nolinkurl{verify.py} re-creates it,
checks all 262,143 nonempty exponent words of total sum at most 18, and
checks every canonical shortened parity source through length 14. These
word tests are regression evidence; the unrestricted quantifier comes
from Sections 2-4.

The crossing audit is a DIFFERENT complete bounded proof domain with
unbounded source and last-exponent parameters. Both ordinary and
optimized runs reproduce verification.json byte for byte. All checks use
explicit exceptions rather than assert statements, so optimization
cannot remove them. Invalid-input controls include a false descent at 1,
a false endpoint descent at 165, an altered source lattice, a false
valuation and a divided nonintegral first-jet witness.

No Lean verification, independent external mathematical review or new
remote GitHub execution is implied by these files. The delivery receipt
records the actual publication state. Earlier valid source files and
their original scope statements are not rewritten by this additive
continuation.

\section{Sources}

{[}MAT{]} A. Languasco, F. Luca, P. Moree and A. Togbe, \emph{Sequences
of integers generated by two fixed primes}, Abh. Math. Semin. Univ.
Hambg. 95 (2025), 123-148, Theorem 2.1. The inspected rational-number
bound is attributed there to Matveev, in a version due to Bugeaud,
Mignotte and Siksek. The proof here uses only that displayed bound and
checks its constants explicitly.
\nolinkurl{https://doi.org/10.1007/s12188-025-00293-9}

{[}RT{]} O. Rozier and C. Terracol, \emph{Paradoxical behavior in
Collatz sequences}, arXiv:2502.00948v2 (13 February 2025), Definitions
1.1-1.2, Sections 2 and 4, and the explicit comparison of iteration
conventions in Section 5. This supplies context and attribution, not a
convergence premise or numerical bound imported into our proof.
\nolinkurl{https://arxiv.org/html/2502.00948v2}

{[}WB1{]} KokunoYumeto, Collatz workbench,
\nolinkurl{intrinsic_zero_firstjet_20260915/firstjet.py}, Git blob
\nolinkurl{97f9ef5ad535f7ec18259d41fab6a3557be17d98}, and its complete
\nolinkurl{note.md} at
\nolinkurl{9afcbe2dc9bace77c1ae705b1cd128da22fe8f8b}. Original integral
first-jet presentation and finite-witness extraction; used unchanged.

{[}WB2{]} Same repository and revision,
\nolinkurl{mixed_switch_control_20260916/note.md} and
\nolinkurl{TASK_STATE.md}. The source stages the
nonrepeating/general-switch target. Its previous mixed-core gap
certificate is not a premise of the new all-word gap.

{[}SZ{]} Owner's Split-Zero construction and intrinsic fibre/quotient
programme, source chain recorded in {[}WB1{]}. The precise structure
used is its comparison kernel with original representatives, relations
and support inclusions, not a claim that an isolated scalar zero
contains a unique trajectory.
